%!TEX root = ../thesis.tex

\chapter{Failure Modes and Diagnosis}\label{chap:failure-modes}

% Claimed as a contribution (\cref{sec:introduction-contributions}): a taxonomy of how
% model-based RL fails on fuzzing-shaped reward landscapes, and the instrumentation
% that tells the modes apart. Frame it that way throughout --- as a methodology for
% diagnosis, evidenced by the failures encountered, not as an account of debugging.

\section{Overview}\label{sec:failure-overview}

% TODO: Set up the chapter: on these targets a run does not degrade gracefully, it
% collapses --- which is why \cref{sec:methodology-seeds} treats the outcome as
% bimodal. Give the taxonomy as a table: mode, mechanism, the diagnostic signature
% that identifies it, and the mitigation. Then one section per mode, each following
% that same order: symptom, mechanism, evidence, mitigation, residual risk.

\section{Policy Collapse and the Echo Chamber}\label{sec:failure-echo-chamber}

% TODO: The prior narrows, search then samples from the narrowed prior, and the
% resulting targets reinforce it. Distinguish it from healthy convergence --- that
% distinction is the hard part and the reason the diagnostic exists. Evidence: the
% per-position argmax matrix over training, MCTS visit-count entropy per depth, and
% the divergence between the search policy and the network prior. If search merely
% echoes the prior, it has stopped improving it.

\section{Dead-Branch Value Fantasy}\label{sec:failure-dead-branch}

% TODO: The model assigns value to branches that terminate, so the planner walks into
% them; the mechanism interacts with how termination and absorbing states are handled
% and with the episode horizon. Evidence: model introspection along a fixed action
% path, and value calibration against realized return. Mitigation and its cost, plus
% the connection to the horizon bound argued in
% \cref{sec:input-reward,sec:corpus-reward}.

\section{Value-Prefix and Support Pathologies}\label{sec:failure-value-prefix}

% TODO: Failures traced to the categorical value representation rather than to the
% policy: a support whose bins are too coarse for the reward gap makes the heads
% action-blind, and an inverted or saturated value prefix silently destroys the
% multi-step signal. This is the empirical justification for deriving the support from
% the reward structure in \cref{sec:methodology-hyperparameters} --- the derivation is
% not a formality, it is a fix for an observed failure.

\section{Reward Hacking}\label{sec:failure-reward-hacking}

% TODO: Coverage as a reward is game-able, and it was gamed: process teardown emitted
% a coverage burst, so dying became profitable. Present it as a first-class finding
% about coverage-as-reward rather than an implementation bug --- any coverage-rewarded
% agent can find teardown paths, and this is the concrete demonstration. Give the fix
% (gate novelty on survival, structure the terminal payout) and then the honest
% caveat: the fix is target-specific reward engineering, which is precisely the
% external-validity problem of \cref{sec:discussion-threats}.

\section{Diagnosing a Collapsed Run}\label{sec:failure-procedure}

% TODO: The transferable payload of this chapter: a procedure. Given a collapsed run,
% which probe to read first, which signature distinguishes which mode, and what to
% change. State which modes the procedure cannot yet distinguish --- the honest limit
% of the method is part of the contribution.
